%==============================================================================
% List of Abbreviations
%==============================================================================
% FIX (bidi + longtable interaction): when a longtable starts on the same
% page as this \chapter*, the title line's node order gets reversed
% (TeX--XeT RTL line), which lands the title flush LEFT. Wrapping the title
% in \hbox to \hsize{...\hss} restores flush-RIGHT geometry under that
% state (verified pixel-wise: title x560-710 like the other front-matter
% chapter* pages). Keep this wrap together with the longtable below —
% without the table the wrap would flip the title left.
\chapter*{\hbox to \hsize{قائمة المختصرات\hss}}
\addcontentsline{toc}{chapter}{قائمة المختصرات}
\markboth{قائمة المختصرات}{قائمة المختصرات}

\renewcommand{\arraystretch}{1.5}
\begin{longtable}{|>{\raggedleft\arraybackslash}p{4cm}|>{\raggedright\arraybackslash}p{11cm}|}
\hline
\textbf{الاختصار} & \textbf{التعريف} \\
\hline
\endhead

\EN{API}  & واجهة برمجة التطبيقات التي تتيح التواصل بين الأنظمة المختلفة. \\
\hline
\EN{BIRD} & معيار قياسي لتقييم أنظمة \EN{NL2SQL} على قواعد بيانات واقعية. \\
\hline
\EN{CSV}  & صيغة ملف نصي لتخزين البيانات الجدولية مفصولة بفواصل. \\
\hline
\EN{ERD}  & مخطط العلاقات بين الكيانات في قاعدة البيانات (\EN{Entity-Relationship Diagram}). \\
\hline
\EN{GPU}  & معالج الرسوميات المستخدم لتسريع معالجة النماذج. \\
\hline
\EN{JWT}  & رموز الويب المميّزة المستخدمة لمصادقة الجلسات وتوقيع البيانات المتبادلة بين الخدمات (\EN{JSON Web Token}). \\
\hline
\EN{LangGraph} & مكتبة لبناء وتنسيق وكلاء الذكاء الاصطناعي المتعددين. \\
\hline
\EN{LLM}  & النموذج اللغوي الكبير (\EN{Large Language Model}). \\
\hline
\EN{MoSCoW} & منهجية ترتيب أولويات المتطلبات: يجب/ينبغي/يمكن/لن (\EN{Must, Should, Could, Won't}). \\
\hline
\EN{NL2SQL} & تحويل اللغة الطبيعية إلى استعلامات \EN{SQL} (\EN{Natural Language to SQL}). \\
\hline
\EN{Odoo} & نظام \EN{ERP} مفتوح المصدر، يُستخدم كبيئة محاكاة جامعية. \\
\hline
\EN{Onboarding} & عملية الإعداد والتهيئة الأولية للنظام. \\
\hline
\EN{pgvector} & امتداد \EN{PostgreSQL} لدعم تخزين والبحث في المتجهات. \\
\hline
\EN{Qdrant} & قاعدة بيانات متجهة مفتوحة المصدر. \\
\hline
\EN{RAG}  & تقنية الاسترجاع المعزز بالتوليد (البحث في المستندات ثم توليد إجابة اعتمادًا على نتائج البحث). \\
\hline
\EN{RBAC} & التحكم في الوصول المبني على الأدوار (\EN{Role-Based Access Control}). \\
\hline
\EN{RTL}  & الكتابة من اليمين إلى اليسار — دعم اللغة العربية. \\
\hline
\EN{SaaS} & البرمجية كخدمة تُقدَّم عبر الإنترنت دون الحاجة لتثبيتها محليًا. \\
\hline
\EN{Schema Linking} & تقنية تصفية الجداول والأعمدة ذات الصلة قبل توليد \EN{SQL}. \\
\hline
\EN{Spider} & معيار قياسي لتقييم أنظمة \EN{NL2SQL}. \\
\hline
\EN{SQL}  & لغة الاستعلام البنيوية المستخدمة للتعامل مع قواعد البيانات. \\
\hline
\EN{SUS}  & مقياس قابلية الاستخدام النظامي (\EN{System Usability Scale}). \\
\hline
\EN{TLS}  & بروتوكول أمن طبقة النقل لتشفير البيانات المنقولة عبر الشبكة (\EN{Transport Layer Security}). \\
\hline
\EN{UAT}  & اختبار قبول المستخدم قبل الإطلاق النهائي (\EN{User Acceptance Testing}). \\
\hline
\EN{UML}  & لغة النمذجة الموحدة (\EN{Unified Modeling Language}). \\
\hline
\EN{Views} & طبقة وسيطة مبسطة فوق قاعدة البيانات لرفع دقة \EN{NL2SQL}. \\
\hline

\end{longtable}
\clearpage
